
\documentclass[11pt]{article}
\renewcommand{\baselinestretch}{1} 
\usepackage[
  a4paper,
  margin=0.75in,
  headsep=3pt, % separation between header rule and text
]{geometry}
\usepackage[most]{tcolorbox}
\setlength{\floatsep}{3pt plus 1pt minus 1pt}
\setlength{\textfloatsep}{3pt plus 1pt minus 1pt}
\setlength{\intextsep}{3pt plus 1pt minus 1pt}
\usepackage{xcolor}
\usepackage{fancyhdr}
\usepackage{scholax}
\usepackage{lastpage}
\usepackage{lineno}
% \linenumbers
\usepackage[inline]{enumitem}
\usepackage{caption}
\usepackage{wrapfig}
\usepackage{graphicx}
\usepackage{ragged2e}
\usepackage{blindtext}
\usepackage{subcaption}
\usepackage{xcolor}
\usepackage{marginnote}
\usepackage{caption}
\usepackage{enumitem}
\setlist[itemize]{noitemsep}
\newcommand{\tcolor}{black}
\DeclareCaptionFont{customsizefig}{\fontsize{10pt}{10pt}\selectfont}
\DeclareCaptionFont{customsizesubfig}{\fontsize{9pt}{9pt}\selectfont}
\captionsetup[figure]{font=customsizefig,aboveskip=1.5pt, belowskip=1.5pt}
\captionsetup[subfigure]{font=customsizesubfig,aboveskip=1.5pt, belowskip=1.5pt}
\newcommand{\TODO}[1]{%
  \marginnote{\color{red}\textbf{TODO}}% Put 'TODO' in margin
  {\color{red}#1}% Put the associated text in red in the body
}
\setlist[itemize]{itemsep=-0.1cm, topsep=0pt, leftmargin=1.2em}
%% ---------------------------------------------------------------
%% Bibliography: natbib + BibTeX, author-year style.  Cite in-line
%% with \citep{key} (parenthetical) or \citet{key} (textual).
%% ---------------------------------------------------------------
\usepackage[round,authoryear]{natbib}
\bibliographystyle{ieeetr}
\setlength{\bibsep}{2pt plus 0.3ex}

\pagestyle{fancy}
\fancyhf{}
\fancyhead[C]{%
  \footnotesize\rmfamily
  Web: \textcolor{blue}{\yourweb}\quad
  \yourname\ \quad
  E-mail: \textcolor{blue}{\youremail}}
\fancyfoot[C]{Page \thepage\ of \pageref{LastPage}}

%% ---------------------------------------------------------------
%% Personal-info macros: edit these once, they propagate everywhere.
%% ---------------------------------------------------------------
\newcommand{\yourname}{Sean Patrick MacBride}
\newcommand{\youremail}{\href{mailto:sean.macbride@physik.uzh.ch}{sean.macbride@physik.uzh.ch}}
\newcommand{\yourweb}{\href{https://seanmacb.github.io}{seanmacb.github.io}}

\newcommand{\hindex}{8}
\newcommand{\hindexDate}{August 2026}

\newcommand{\soptitle}{Research Statement}

%% Section-header macro (kept from the original SOP template)
\newcommand{\statement}[1]{\par\medskip
  \textcolor{blue}{\textbf{#1:}}%\hspace{0.05cm}
}

%% Highlight macro for spots the author should tailor per application
%% (e.g. naming the specific host group, fellowship, or institution).
\newcommand{\fillin}[1]{\textcolor{orange}{\textbf{[#1]}}}

%% ---------------------------------------------------------------
%% Recurring survey / collaboration / instrument acronyms.
%%
%% These are now managed by the `acro' package: each is declared once
%% with a long and a short form, and \ac{...} automatically prints
%% "Long Form (SHORT)" on first use and "SHORT" every time after.
%% The user-facing macros (\LSST, \DESC, ...) are kept unchanged so
%% the body text below did not have to be rewritten.
%% ---------------------------------------------------------------
\usepackage{acro}
\acsetup{make-links=false}

\DeclareAcronym{lsst}    {short = LSST,    long = Legacy Survey of Space and Time}
\DeclareAcronym{lsstcam} {short = LSSTCam, long = LSST Camera}
\DeclareAcronym{gw} {short = GW, long = Gravitational wave}
\DeclareAcronym{em} {short = EM, long = electromagnetic}
\DeclareAcronym{desc}    {short = DESC,    long = LSST-Dark Energy Science Collaboration}
\DeclareAcronym{des}     {short = DES,     long = Dark Energy Survey}
\DeclareAcronym{desi}    {short = DESI,    long = \textbf{Dark Energy Spectroscopic Instrument}}
\DeclareAcronym{desgw}   {short = DESGW,   long = Dark Energy Survey Gravitational Wave follow-up program}
\DeclareAcronym{delve}   {short = DELVE,   long = Dark Energy Local Volume Exploration}
\DeclareAcronym{pfs}     {short = PFS,     long = Prime Focus Spectrograph}
\DeclareAcronym{llamas}  {short = LLAMAS,  long = \textbf{Large Lenslet Array Magellan Spectrograph}}
\DeclareAcronym{lvk}     {short = LVK,     long = LIGO-Virgo-Kagra}
\DeclareAcronym{too}     {short = ToO,     long = Target-of-Opportunity}
\DeclareAcronym{decam}   {short = DECam,   long = \textbf{Dark Energy Camera}}
\DeclareAcronym{rubin}   {short = Rubin,   long = Vera C. Rubin Observatory}
\DeclareAcronym{must}   {short = MUST,   long = MUltiplexed Survey Telescope}
\DeclareAcronym{3G}   {short = 3G,   long = third generation}
\DeclareAcronym{specsfive}   {short = Spec-S5,   long = \textbf{Stage-5 Spectroscopic Experiment}}
\DeclareAcronym{LISA}   {short = LISA,   long = Laser Interferometer Space Antenna}
\DeclareAcronym{hyperk}   {short = Hyper-K,   long = Hyper Kamiokande}
\DeclareAcronym{DUNE}   {short = DUNE,   long = Deep Underground Neutrino Experiment}
\DeclareAcronym{jgem}    {short = J-GEM,
                          long  = Japanese collaboration for gravitational-wave electromagnetic follow-up}
\DeclareAcronym{UNIONS}   {short = UNIONS,   long = The Ultraviolet Near-infrared Optical Northern Survey} 

\newcommand{\LSST}{\ac{lsst}}
\newcommand{\EM}{\ac{em}}
\newcommand{\threeG}{\ac{3G}}
\newcommand{\UNIONS}{\ac{UNIONS}}
\newcommand{\LSSTCam}{\ac{lsstcam}}
\newcommand{\DESC}{\ac{desc}}
\newcommand{\DES}{\ac{des}}
\newcommand{\DESI}{\ac{desi}}
\newcommand{\DESGW}{\ac{desgw}}
\newcommand{\DELVE}{\ac{delve}}
\newcommand{\PFS}{\ac{pfs}}
\newcommand{\LLAMAS}{\ac{llamas}}
\newcommand{\LVK}{\ac{lvk}}
\newcommand{\ToO}{\ac{too}}
\newcommand{\DECam}{\ac{decam}}
\newcommand{\JGEM}{\ac{jgem}}
\newcommand{\GW}{\ac{gw}}
\newcommand{\must}{\ac{must}}

\newcommand{\specsfive}{\ac{specsfive}}
\newcommand{\LISA}{\ac{LISA}}
\newcommand{\hyperk}{\ac{hyperk}}
\newcommand{\DUNE}{\ac{DUNE}}

%% Non-acronym text macros (unchanged: these are names, not abbreviations)
\newcommand{\Rubin}{\ac{rubin}}
\newcommand{\Ho}{\ensuremath{H_0}}
\newcommand{\Blanco}{V\'ictor M. Blanco 4-meter Telescope}
\newcommand{\IceCube}{IceCube Neutrino Observatory}

\newcommand{\Institute}{\textcolor{\tcolor}{Stanford University}}
\newcommand{\Position}{\textcolor{\tcolor}{Stanford Science Fellow}}

\usepackage[
  breaklinks,
  pdftitle={\yourname - \soptitle},
  pdfauthor={\yourname},
  unicode
]{hyperref}

\begin{document}
\hrule
\begin{center}
    \Large\soptitle\ -- \yourname
\end{center}
\vspace{0pt}
\hrule
\vspace{0.5pt}
\hrule height 0.5pt
\vspace{3pt}

\normalsize
\thispagestyle{plain}
\statement{Cracks in Modern Cosmology}
The standard model of cosmology, $\Lambda$CDM, provides a precise description of the universe and many of its defining features. As observational cosmology has entered a golden age of high-precision, some parameters have shown discrepancies that warrant extensive study and assessment, including the local rate of accelerated expansion, the Hubble constant ($H_0$). $H_0$ measurements derived from the early Universe -- most precisely, the \textit{Planck} cosmic microwave background analysis, infer $H_0 = 67.4\pm0.5\,\mathrm{km\,s^{-1}\,Mpc^{-1}}$ using the temperature power spectrum, E-mode polarization, and cross correlations \citep{Planck2020} - disagreeing with distance-ladder measurements from the local universe, most notably the SH0ES result of $H_0 = 73.0\pm1.0\,\mathrm{km\,s^{-1}\,Mpc^{-1}}$ \citep{Riess2022}. 
Despite extensive scrutiny of each approach, no systematic effect or new physics has been identified to resolve the discrepancy.
This $>5\sigma$ discrepancy is known as the Hubble Tension. 

\GW\ standard sirens offer an independent measurement of distances, and therefore, $H_0$. The \GW\ waveform from the coalescence of two compact objects directly encodes the luminosity distance to the source \citep{Schutz1986}. However, a \GW-based measurement alone cannot provide the source redshift, and several strategies have been developed to measure cosmological parameters using the \GW\ waveform information. Two strategies have been developed to combine \GW\ and electromagnetic data to enable a standard siren measurement, and are the focus of my research.

The first, the \textit{bright siren method}, requires the detection of an \EM\ counterpart to a \GW\ merger. The identification of a \GW\ merger candidate and the spectroscopic redshift of its host galaxy provides the highest $H_0$ precision per \GW\ event. The binary neutron star merger GW170817 remains the only \GW\ event with a confirmed \EM\ counterpart to date, and provided the first standard-siren measurement of $H_0$ \citep{2017ApJ...848L..12A,Abbott2017Nature,SoaresSantos2017}. The second method is the \textit{dark siren method}. Using dark sirens, it is possible to statistically marginalize over galaxies consistent with a \GW\ event's three-dimensional localization volume, weighting each galaxy by its probability of hosting the merger \citep{Gray2023}. This approach requires no \EM\ counterpart, and can be applied to all \GW\ detections coincident with galaxy catalog data.

I have designed a research program that will leverage standard sirens as an emerging probe of cosmology, maximizing the data from \LSST\ for bright and dark sirens, developing other astronomical instruments to leverage \LSST\ data, and ultimately using gravitational wave and electromagnetic data to address the Hubble Tension. 
\statement{Standard Sirens in the LSST Era}
My research is grounded in standard siren cosmology, but spans astronomical instrument development, bright siren searches, and applications of the dark siren method to \LSST\ data.

\textbf{Astronomical Instruments to Leverage LSST data:}
In the \LSST\ era, the central bottleneck will not be imaging depth, but rather spectroscopic capacity to characterize \LSST\ transients. To aid in this effort, I will design and deploy a new multi-object spectrograph optimized for standard siren cosmology, tuning the focal plane parameters for the precision and event rates of forthcoming gravitational wave detectors. This instrument would serve two goals: 1) characterize astrophysical transients associated with bright standard sirens, and 2) obtain spectroscopic redshifts from host galaxies to bolster dark siren measurements. Through this dual purpose, the proposed instrument would become a dark energy science machine, primed to capitalize on the data deluge from \LSST. In addition to \GW\ follow-up, this instrument could characterize any transient from the \LSST\ alert stream, augmenting supernovae cosmology measurement samples. Given the scale of this instrument, this instrument could be completed and commissioned within five years.
% These pathfinding studies will inform the proposed \specsfive, a future, large-scale cosmic frontier experiment planned for the mid-2030's \cite{2025arXiv250307923B}. Developing focal-plane technology for other instruments will help shape the capabilities of \specsfive, so it serves as an effective spectroscopic complement to \Rubin. % Spec-S5 instrumentation development.

% \textbf{Instrumentation Development for Future Cosmic Surveys:}
% Building on my experience commissioning and testing CCDs for the LSST Camera focal plane, I propose to characterize the next generation of astronomical sensors at Fermilab, focusing on the multi-amplifier sensing (MAS) CCDs slated for \textbf{\DESI\ Run 2} and the \specsfive. This work will apply my expertise \LSSTCam\ characterization to test, qualify, and optimize MAS CCDs for forthcoming spectroscopic surveys. In doing so, I will ensure that the detector performance meets the underlying requirements of \DESI\ Run 2 and \specsfive, maintaining the precision cosmology goals from the project's inception, while contributing directly to Fermilab's detector development program. This testing campaign will have impacts in the near-term, as \DESI\ Run 2 is expected to start in 2029, along with pathfinding for a longer time-horizon experiment, \specsfive.
In addition to designing and building this future instrument, 
I will sustain a support-scientist role for LSSTCam as the \LSST\ transitions into survey operations, maintaining the operational knowledge needed to respond quickly to sensor anomalies and focal-plane issues as they arise.
Building on my commissioning-era characterization of \LSSTCam\ defects, I will also deploy a defect-monitoring tool, currently in-development, to \texttt{cp\_verify}, the calibration-product verification framework used by the \LSST\ Science Pipelines \cite{10.71929/rubin/2570545}. This tool will elevate defect monitoring to an automated process, improving \LSST\ data quality for all science areas.

\textbf{Bright Siren Searches Enabled by Rubin Observatory:}
I will lead the observing planning, execution, and analysis of the next bright standard siren accessible to \Rubin\ Observatory, 
% \begin{wrapfigure}{r}{0.55\textwidth} % 'r' for right, 'l' for left
%     \centering
%     \includegraphics[width=0.5\textwidth]{figures/pfs_completeness.png}
% 	\caption{In red, the completeness of \PFS\ objects acquired as part of our \ToO\ campaign \cite{Zhang2025S250328ae}, with other colors owing the forecasted completeness of proposed observations. The high completeness obtained over the localization region in 3 passes (one night) demonstrates the utility of current and future multi-object spectrographs as characterization machines for standard siren cosmology.}
%     \label{fig:pfs_coverage}
% \end{wrapfigure}
along with other \ToO's pursued during the \LSST, continuing my role as the \Rubin\ \ToO\ observer throughout the transition to operations. With current timelines, the next opportunity to search for \EM\ counterparts to \GW s is during the \LVK\ Observing Run 5 (O5), in 2029, after the intermediate run that ends in mid 2027.
% This leadership is enabled by my specialized expertise in the alert processing of the \ToO\ system, my intimate knowledge of the technical strengths and limitations of the observatory, and my sustained development of the \ToO\ system within the wider \Rubin\ ecosystem.
While it is difficult to predict the timing of the next bright siren, I will continue to deploy improvements to the \ToO\ system to improve response latency, observing strategy, and automation, aiming to maximize the next opportunity to detect and characterize these 
optical counterparts and obtain a well-constrained measure of \Ho.
% Lead DESGW IRX analyses - observation planning, image processing, and transient analysis - drafted

In parallel to operating the \Rubin\ \ToO\ system, I will analyze transient candidates from the \DESGW\ program in future \LVK\ observing runs. This positions me to sustain and support \Rubin's \ToO\ program with additional imaging in other filters and extend the temporal-sampling using a similarly capable photometric imager. This complimentary support will play an important role, especially as other priorities of the \LSST\ are balanced against the \ToO\ program.
% Other observing programs with spectroscopy to compliment the Rubin alert stream - in progress
As the \LSST\ \ToO\ program matures, I will work to enhance transient characterization by securing spectroscopic follow-up time on other instruments, moving beyond photometric classification to confirm spectral features that can inform the scope of \Rubin's \ToO s, and constrain the physical properties of \LSST\ transients.% and their host environments. 

%% LSST efforts for dark sirens
\textbf{Dark Siren Measurements for Precision Cosmology:}
I will pursue a cutting-edge program that creates a gold-standard photometric galaxy catalog for dark siren cosmology using \LSST\ data.
% LSST DP2 characterization in prep for DR1
This program is structured in four stages, each taking approximately one year, and building towards a powerful, world leading measurement of $H_0$ using \GW\ and \EM\ data.
It begins with an ongoing project, focused on characterizing the \LSST\ Data Preview 2 (DP2) galaxy catalog, aiming to define robust and scalable object classification for future \LSST\ data releases%\cite{sourceClass}
.
% DESC Y1 forecasting
Following characterization of DP2, I will finish a forecasting study that quantifies the $H_0$ precision achievable by combining simulated first-year LSST data with an O4-like gravitational-wave catalog, extending the dark-siren methodology from my recent \DELVE\ measurement \citep{MacBride2026DELVEH0} to project the constraining power in the \LSST\ era.
% LSST DR1 measurement
Superceding DP2 preparation, I will characterize the first \LSST\ Data Release (DR1) galaxy catalog and produce a dark siren measurement of $H_0$. This result will extend the methodology I developed with \DELVE\ to a dataset with substantially greater depth, area, and photometric uniformity, all features that will make it the standard galaxy catalog for dark siren analyses.
% LSST DR1 combined probes - emphasize shear instead of clustering
Building on the DR1 measurement, I aim to combine \LSST\ large-scale structure measurements in a joint cosmological analysis, using cosmic shear and standard sirens to obtain a tighter constraint of $H_0$, and extend the constraints to $\Omega_m$.

%% Adapting other data for Dark sirens, either through specialized surveys, or through existing data
% GWCats - mention when discussing the existing collaboration for GW-Primo
As the wealth of \GW\ data increases alongside a wide variety of new galaxy catalog data, I will explore the frontier of the standard siren method using specialized surveys.
%and a rigorous assessment of systematics for standard sirens,
% while engaging with the Horizon Europe-funded GWCats collaboration, a tight-knit group of researchers focused on using galaxy catalogs and \GW\ data to constrain cosmology.
% NORTHSTARS
This involves leading a dedicated spectroscopic dark-siren pilot survey, obtaining targeted redshifts within GW localization volumes to assess the limits and enhancements to $H_0$ constraints using well-localized dark sirens. This survey is led through a continuing collaboration with \JGEM\ on the \PFS\ instrument, obtaining additional observations for the GW merger pursued in our collaborative \DESGW\ \ToO, with an expected observation date starting in April 2027 and analysis published by the end of 2027. After completing the pilot survey, we will expand the program to other \GW\ events.
% GW-primo
% In parallel, I will lead the image processing for a similar survey, the Gravitational Wave Photometric Redshift Improvement with Medium-band Observations (GW-Primo) project, to test how medium band observations improve $H_0$ precision through enhanced photometric redshifts. By comparing the results of these two survey programs, I will be able to assess which route, enhanced photometric or spectroscopic, will provide the best constraints on $H_0$ given limited telescope time and the current sample of well-localized GWs.

% UNIONS
% Similar to \DES\, the \UNIONS\ survey provides a large area galaxy catalog in the Northern hemisphere, where \DES\ and \LSST\ coverage is not possible. I will prepare the  \UNIONS\ galaxy catalog for dark siren cosmology through ongoing work with the \UNIONS\ collaboration, applying the lessons from the \DELVE\ galaxy catalog preparation. I will construct both an isolated \UNIONS\ catalog and a combined galaxy catalog of \UNIONS\ and \LSST\ data for standard siren cosmology, while taking care to control any systematics that come from joining different galaxy catalogs. This combined catalog will span most of the accessible extragalactic sky, including several well localized GWs in the northern hemisphere, becoming a value-added catalog compared to \LSST\ data alone.

% PV methodological study
% Finally, with the new discovery capabilities of \Rubin\ for bright standard sirens, the systematics of multiple bright siren measurements must be thoroughly examined before claiming a robus and precise measurement of $H_0$. To build this foundation, I will lead an assessment of peculiar velocity systematics in bright siren measurements, using synthetic and real data to accurately map the velocity field and reduce the systematic contribution to future bright siren measurements.

\statement{How Stanford is Essential}
Stanford and SLAC bring together every capability this program requires. 
SLAC built \LSSTCam, and hosts the U.S. Data Facility, placing my camera-support, \texttt{cp\_verify} defect-monitoring, and \ToO\ work alongside the engineers and pipeline developers who deploy it. 
KIPAC's leadership in the \LSST\ Dark Energy Science Collaboration is the ideal setting to extend the DR1 dark siren measurement into a joint shear and standard siren constraint on $H_0$ and $\Omega_m$. 
SLAC's expertise in delivering large-scale instruments, including \LSSTCam, together with Stanford's observational cosmology community, provides the engineering depth needed to support the design of the proposed multi-object spectrograph. 
The combination of these strengths will elevate this research program to develop the gold-standard measurement using standard sirens, directly addressing the Hubble Tension.

\newpage

\textbf{References}
% {\small A full publication list is available on \href{https://scholar.google.com/citations?user=XtRTswUAAAAJ}{Google Scholar} and \href{https://ui.adsabs.harvard.edu/search/q=((author%3A%22MacBride%2C%20S.%22)%20AND%20year%3A2017-2100)&sort=date%20desc%2C%20bibcode%20desc&ui_tag=results%2Fgraph%2Fyear&p_=0}{NASA ADS} (h-index \hindex\ as of \hindexDate).}
\renewcommand{\refname}{}
\vspace{-3\baselineskip}
{\small\bibliography{references}}

\end{document}